\chapter{Antecedentes y Alcance}

\section{Contexto del Proyecto}

El sistema eléctrico colombiano ha dependido históricamente de la generación hidráulica, que concentra entre el 70\% y el 80\% de la energía producida. Si bien esa matriz ha resultado competitiva en costos, también expone al país a riesgos asociados a la variabilidad hidrológica y a eventos climáticos extremos. En respuesta, el Estado ha impulsado la diversificación mediante Fuentes No Convencionales de Energía Renovable (FNCER), creando un marco legal y regulatorio que facilita la incorporación de generación distribuida en los sistemas de distribución.

Entre los instrumentos más relevantes se encuentran la Ley~1715 de 2014 y sus actualizaciones; la Resolución CREG~030 de 2018, que regula la conexión de autogeneradores a pequeña y gran escala (AGPE/AGGE); la Resolución CREG~174 de 2021, que precisa lineamientos de autogeneración a pequeña escala; y la Resolución CREG~025 de 1995 ---Código de Redes---, con sus modificaciones posteriores. Este marco define tanto las obligaciones del interesado como los criterios técnicos que el operador de red debe verificar antes de autorizar una conexión.

\section{Alcance del Estudio}

En ese contexto, el supermercado MAS~X~MENOS sede Guarín adelanta el proyecto \textbf{SOLAR MAS X MENOS GUARÍN}, con una potencia instalada de 141,6~kWp. Dado que la potencia en corriente alterna es de 120~kW, la instalación se clasifica como autogenerador a pequeña escala (AGPE) y, por tanto, requiere un estudio de conexión simplificado ante ESSA.

El alcance técnico del presente trabajo cubre la caracterización del circuito 10-502 (SE~Conucos) a partir de los datos oficiales del operador; el modelado del punto de conexión en el nodo 3272966 y del transformador de 150~kVA; el análisis de flujo de carga en operación normal para las franjas de las 9:00, 12:00 y 15:00; el análisis de cortocircuito trifásico y monofásico conforme a la norma IEC~60909; la verificación de las protecciones del circuito y del generador a la luz del Acuerdo CNO~1322; y, finalmente, la estimación de las pérdidas técnicas con y sin el sistema solar fotovoltaico (SSFV).

\section{Objetivo General}

El objetivo general consiste en analizar la viabilidad técnica de conectar el proyecto fotovoltaico de 141,6~kWp en el nodo 3272966 del circuito 10-502 de ESSA (13,8~kV), cuantificando sus efectos sobre tensiones, cargabilidades, niveles de cortocircuito, esquema de protecciones y pérdidas del alimentador, de modo que el operador cuente con un sustento técnico claro para decidir sobre la solicitud de conexión.

\section{Objetivos Específicos}

Para alcanzar ese propósito se plantean objetivos específicos que orientan cada etapa del estudio. En primer lugar, se caracteriza la red de influencia utilizando los parámetros eléctricos de 215 tramos y 51 cargas del circuito. A continuación, se evalúan las tensiones y cargabilidades con y sin generación en las franjas horarias críticas. Se calculan, además, las corrientes de cortocircuito trifásico y monofásico en el punto de conexión y en nodos aledaños. El estudio también contrasta los ajustes del relé MICOM~P142 del circuito 10-502 frente al aporte del SSFV y desarrolla la coordinación de protecciones mediante curvas tiempo--corriente. Por último, se cuantifica la variación de pérdidas técnicas atribuible a la generación distribuida y se emite el concepto técnico de conexión, junto con las recomendaciones de seguimiento pertinentes.
